Table~\ref{tab:rerank} applies the exposure-capped allocation of
Theorem~\ref{thm:alloc} to every operator. The min-cost-flow solver returns the
optimum up to its rounding gap (a relative gap below $4\cdot10^{-5}$ in every run), all
caps hold, and coverage never falls below the guaranteed $NK/\max\mathrm{cap}$; a
dual-subgradient solver we tried first left gaps of up to 35\%, which the
certificate exposed. On KuaiRec, where exposure is extremely concentrated,
DR adjacency and DRUP are the most accurate operators under every cap, for
example 0.199 against 0.173 (logged graph) and 0.178 (IPS) at $c=5$, and 0.140
against 0.115 and 0.121 at $c=2$. Under the same cap, however, their lists
remain more concentrated (Gini 0.80 against 0.72 at $c=5$), and at equal Gini
the frontier over each operator's grid (Table~\ref{tab:frontier}) favours the
correlation-based operators on Coat. On Coat, whose random test items are already
diverse, capping costs DRUP more accuracy than it costs the logged graph; on
Yahoo!\,R3 exposure is not concentrated (Gini 0.27) and caps down to twice the
fair share cost every operator at most 0.002. The allocation is thus the
mechanism that controls concentration, and the DR graph is the better input to
it when exposure in the log is very uneven.
